%%%%%%%%%%First Chapter%%%%%%%%%%
\chapter{Introduction}

\section{Background}

% \paragraph{1 - Development of 5G and the Utilization of FR2 Spectrum}
% Paragraf pertama membuka pembahasan dari kebutuhan komunikasi nirkabel berkecepatan tinggi dan penggunaan spektrum frekuensi tinggi pada 5G New Radio.

% Poin utama:

% Pertumbuhan kebutuhan data dan perangkat terhubung membutuhkan bandwidth komunikasi yang lebih luas karena scarcity and fragmentation in conventional cellular bands dan pindah ke mmwave\cite{Rappaport2013}
% Salah satu pengembangan penting dalam 5G New Radio adalah pemanfaatan Frequency Range 2 (FR2) yang memiliki bandwidth yang luas \cite{TSGR2021}.
% Frekuensi penelitian sebesar 38 GHz berada dalam FR2-1.
% Secara lebih spesifik, 38 GHz termasuk dalam NR operating band n260, yang mencakup 37–40 GHz yang menyediakan hampir 3 GHz bandwidth untuk komunikasi nirkabel berkecepatan tinggi.
% % Sitasi : \cite{Rappaport2013,TSGR2021}
The rapid growth in data demand and the proliferation of connected devices necessitate significantly wider communication bandwidths to support next-generation wireless networks. This is primarily driven by the scarcity and fragmentation within conventional cellular bands, pushing the industry toward millimeter-wave (mmWave) frequencies \cite{Rappaport2013}. One of the most significant developments in 5G New Radio (NR) technology is the utilization of high-frequency spectrum, specifically Frequency Range 2 (FR2), which offers extensive bandwidth for high-speed wireless communication \cite{TSGR2021}. This research focuses on a frequency of 38 GHz, which falls within the FR2-1 segment. More specifically, 38 GHz is situated within the NR operating band n260, spanning 37–40 GHz and providing approximately 3 GHz of bandwidth for high-capacity wireless services.

% \paragraph{Propagational Challenges and Hardware Complexity in 5G FR2}
The utilization of the high-frequency FR2 spectrum is inherently linked to significant propagation challenges arising from its short wavelengths. Signals at these frequencies suffer from severe free-space path loss, substantial penetration losses, and extreme sensitivity to physical blockages, which collectively restrict overall coverage \cite{Zeng2016}. Consequently, maintaining reliable wireless links in FR2 environments requires sophisticated techniques such as directional transmissions and high antenna gains. While the small wavelength size facilitates the integration of numerous antenna elements into a compact aperture, it also necessitates advanced spatial processing capabilities. Conventional approaches to achieve this include Massive MIMO systems designed to provide sufficient beamforming gain; however, fully digital MIMO architectures—which require an independent RF chain for every antenna element—pose substantial practical hurdles at FR2 frequencies. The high-frequency environment demands complex and costly components such as power amplifiers (PAs), low-noise amplifiers (LNAs), mixers, local oscillators, and high-speed ADCs or DACs \cite{Heath2016}. Scaling the number of RF chains in these systems significantly escalates hardware complexity, power consumption, manufacturing costs, and the stringency of calibration requirements.

% Poin utama Paragraf 2:
% Sinyal FR2 memiliki panjang gelombang yang sangat pendek.
% Frekuensi tinggi mengalami:
% free-space path loss yang tinggi;
% penetration loss yang besar;
% sensitivitas terhadap blockage;
% cakupan propagasi yang lebih terbatas.
% Karena itu, sistem FR2 memerlukan directional transmission dan antenna gain yang tinggi.
% Ukuran panjang gelombang yang kecil memungkinkan banyak elemen antena ditempatkan pada aperture yang relatif ringkas. \cite{Zeng2016}

% Poin utama Paragraf 3:
% Massive MIMO digunakan untuk menghasilkan beamforming gain.
% Fully digital MIMO membutuhkan satu RF chain untuk setiap elemen antena.
% Pada frekuensi FR2, RF chain memerlukan komponen seperti:
% power amplifier;
% low-noise amplifier;
% mixer;
% local oscillator;
% ADC atau DAC.
% Penambahan jumlah RF chain meningkatkan:
% hardware complexity;
% konsumsi daya;
% biaya implementasi;
% kebutuhan kalibrasi.
% sitasi : \cite{Heath2016}

% \paragraph{Lens Antenna Array for Beamspace Communication in FR2 MIMO}
The lens antenna array presents a compelling alternative to conventional MIMO systems by utilizing an electromagnetic lens in conjunction with antenna elements positioned on the focal region or surface. This configuration allows the lens to provide varying phase delays across the aperture, facilitating the passive spatial transformation of signals from the antenna space into a beamspace \cite{shen2019}. By focusing waves originating from different directions onto specific subsets of antenna elements, only those receiving dominant energy need to be connected to the RF chains. This characteristic is particularly advantageous for Frequency Range 2 (FR2) MIMO systems, where propagation channels are typically characterized by a limited number of dominant paths and a relatively sparse beamspace \cite{Zeng2016}. In such environments, most channel energy is concentrated on only a few beams or antennas, allowing for efficient beam selection to significantly reduce the required number of RF chains. Furthermore, the integration of lens antennas can enhance system performance by improving angular separation, spatial channel decorrelation, effective channel rank, spatial multiplexing capability, and overall channel capacity \cite{Zeng2016}.

% Poin utama Paragraf 4:
% Lens antenna array terdiri dari:
% electromagnetic lens;
% antenna elements pada focal region atau focal surface.
% Lens memberikan phase delay berbeda di sepanjang aperture.
% Gelombang dari arah berbeda difokuskan pada elemen antena berbeda.
% Lens secara pasif melakukan spatial transformation dari antenna space menuju beamspace.
% Hanya antena yang menerima energi dominan yang perlu dihubungkan ke RF chain.

% Poin utama Paragraf 5:
% Kanal FR2 umumnya memiliki jumlah dominant propagation paths yang terbatas.
% Beamspace channel menjadi relatif sparse.
% Sebagian besar energi kanal terkonsentrasi pada beberapa beam atau antena.
% Beam selection dapat mengurangi jumlah RF chain.
% Lens juga berpotensi meningkatkan:
% angular separation;
% spatial channel decorrelation;
% effective channel rank;
% spatial multiplexing capability;
% channel capacity.\cite{Zeng2016}

% \paragraph{Metasurface-Based Planar Lens Solutions for High-Frequency Communication}
While conventional dielectric lenses are valued for their high focusing capabilities and gain, they possess several significant drawbacks that hinder their integration into modern wireless systems. These include thick profiles, substantial weight, curved surfaces, and reflection losses stemming from impedance mismatches \cite{Al-Joumayly2011}. Furthermore, dielectric lenses can be bulky, expensive to manufacture at certain microwave frequencies, and inherently difficult to integrate with planar antenna arrays due to their complex geometries. To overcome these limitations, there is a growing shift toward planar lens architectures. A prominent solution in this area is the metasurface-based planar antenna lens, which comprises an array of sub-wavelength unit cells that act as individual spatial phase shifters \cite{Al-Joumayly2011}. By adjusting the geometric dimensions of these unit cells, one can precisely control the transmission magnitude, phase, and frequency response across the aperture. The distribution of these cells allows for the synthesis of specific phase profiles capable of transforming spherical wavefronts into planar ones or vice versa, effectively focusing waves while maintaining a compact and low-profile form factor \cite{Al-Joumayly2011}.

% Poin utama Paragraf 6:
% Dielectric lens dapat memberikan focusing dan gain tinggi.
% Namun, dielectric lens dapat memiliki:
% profil tebal;
% bobot besar;
% curved surface;
% reflection loss akibat impedance mismatch;
% integrasi yang sulit dengan planar antenna array;
% proses fabrikasi yang relatif kompleks.
% Masalah ini mendorong pengembangan planar lens.

% Poin utama Paragraf 7:
% Metasurface terdiri dari kumpulan sub-wavelength unit cells.
% Setiap unit cell bertindak sebagai spatial phase shifter.
% Dimensi geometris unit cell menentukan:
% transmission magnitude;
% transmission phase;
% frequency response.
% Distribusi unit cell di seluruh aperture membentuk phase profile yang diperlukan.
% Metasurface dapat mengubah:
% spherical wavefront menjadi planar wavefront;
% planar wavefront menjadi focused wavefront.

% \paragraph{8 — Relevance of a 38 GHz Metasurface Lens to 5G FR2}

% Paragraf ini membuat hubungan langsung antara desainmu dan 5G FR2.

% Poin utama:

% Lens dirancang pada 38 GHz.
% Frekuensi tersebut berada dalam NR band n260.
% Teknologi ini relevan untuk:
% FR2 access points;
% base stations;
% fixed wireless access;
% indoor high-capacity communication;
% directional MIMO receivers.
% Metasurface dapat melakukan passive beam focusing tanpa memerlukan phase shifter aktif pada setiap elemen.


\section{Motivation}
% \paragraph{Design Motivation and Requirements for 5G FR2 Lens Antennas}
The primary motivation for this research stems from the necessity of developing compact and easily integrable devices for 5G FR2 systems. To address the limitations of conventional dielectric lenses, which are often thick, heavy, and curved, there is a significant demand for planar, low-profile lens architectures that can be fabricated using standard PCB techniques. Such structures facilitate seamless integration with patch antennas and focal-plane receivers while maintaining high phase coverage and controlled insertion loss. Furthermore, the proposed metasurface lens leverages passive focusing at 38 GHz to enhance received power without the need for active phase shifters. This design allows the focus to adapt dynamically to the angle of incidence, enabling multiple antenna receivers to be strategically positioned along a focal arc or plane, thereby optimizing spatial coverage and performance in high-frequency environments.

% Poin utama Paragraf 1:
% Mengembangkan lens yang planar dan low-profile.
% Menghindari struktur dielectric lens yang tebal dan melengkung.
% Menggunakan teknik fabrikasi PCB.
% Mengintegrasikan lens dengan patch antenna dan focal-plane receiver.
% Menghasilkan phase coverage yang besar dengan insertion loss yang tetap terkendali.

% Poin utama Paragraf 2:
% Metasurface lens mengontrol phase gelombang tanpa active phase shifter.
% Fokus yang dihasilkan dapat meningkatkan received power.
% Fokus berubah sesuai incidence angle.
% Beberapa antenna receiver dapat ditempatkan pada focal arc atau focal plane.

% \paragraph{Beyond Antenna Gain Enhancement}
It is critical to emphasize that the contributions of this thesis extend beyond mere antenna gain enhancement. While increased gain is a primary objective, the core novelty lies in the lens's ability to fundamentally alter the structure of the communication channel matrix. By facilitating field separation at the focal region, the proposed design effectively reduces channel overlap between spatially distinct paths. Consequently, it is hypothesized that this spatial transformation will lead to a significant reduction in spatial correlation and more balanced singular values within the channel matrix. These improvements are expected to enhance the condition number and increase the effective rank of the channel, ultimately leading to a substantial increase in overall channel capacity for high-frequency MIMO systems.

% Poin utama:
% Lens bukan hanya gain enhancer.
% Lens dapat mengubah struktur channel matrix.
% Pemisahan field pada focal region dapat menurunkan channel overlap.
% Hipotesis penelitian:
% spatial correlation menurun;
% singular values menjadi lebih seimbang;
% condition number membaik;
% effective rank meningkat;
% channel capacity meningkat.

\section{Thesis Structure}
This thesis is organized as follows. Chapter 1 introduces the research background, motivation, and thesis structure. Chapter 2 presents the theoretical background of metasurface transmitarray lens antennas and communication channel evaluation. Chapter 3 describes the simulation setup and design workflow, including unit cell design, transmitarray lens configuration, ANSYS HFSS simulation, and channel evaluation setup. Chapter 4 presents the simulation results and discussion for both antenna-level and communication-level performance. Chapter 5 summarizes the conclusions, main contributions, limitations, and future work.
